% GENERATED FILE. Edit canonical lessons, then run npm run generate:books.
% canonical-source-hash: fnv1a64:52e5f54e8cb113c5
% canonical-lessons: ES-C48-cuando

\chapter{Cuando --- Both Pasts in One Sentence}
\label{ch:cuando}
\hlchaptermodality{107}

\begin{chapteropening}
\textbf{By the end of this chapter:} \emph{I can join a scene and an event with cuando, using the right tense for each.}

Complete ES-C48-cuando: build estudiaba cuando hablo Ana and say which tense sets the room.
\end{chapteropening}

\section[cuando]{cuando — putting both pasts in one sentence}

\label{lesson:ES-C48-cuando}



\begin{quote}
One tense is the scenery and one is the event. So far you have only been able to use one at a time. This chapter's word lets them share a sentence.
\end{quote}

\begin{sounds}
\begin{itemize}
\raggedright
  \item \texttt{diphthong-ua} — \emph{cua} is one syllable, not two: \emph{KWAN-doh}.
\end{itemize}
\end{sounds}

\begin{cousinweb}
\begin{quote}
\textbf{cuando} — \emph{when}.
\end{quote}

From Latin \textbf{quandō}, and the \emph{qu-} at the front is not decoration. It is the same question-marker sitting inside \emph{qué}, \emph{quién} and \emph{cómo} — a family you have been collecting since your very first questions.

Latin built its question words from one root and Spanish kept the whole set, so \emph{cuándo} with an accent asks a question and \emph{cuando} without one simply joins two things. Same word, and the accent is doing the same job it always does: marking the stressed, questioning one.
\end{cousinweb}

\begin{grammarlens}[title={Grammar Lens: the scene and the thing that interrupts it}]
Now the two tenses can do their jobs side by side:

\begin{quote}
\emph{\textbf{Estudiaba} \textbf{cuando} \textbf{habló} Ana.} — I \textbf{was studying} when Ana \textbf{spoke}.
\end{quote}

Read what each tense is doing:

\begin{itemize}
\raggedright
  \item \emph{Estudiaba} — imperfect. The \textbf{scene}. Already going, no edges, and it does not stop being true when something else happens.
  \item \emph{Habló} — preterite. The \textbf{event}. One thing, with edges, dropped into that scene.
\end{itemize}

This is the shape most Spanish stories are built from, and it is worth hearing the pattern rather than memorising a rule: \textbf{the imperfect sets the room, the preterite is what walks into it.}

Turn it around and the sentence changes meaning, not just style:

\begin{quote}
\emph{Estudié cuando hablaba Ana.} — I studied while Ana was talking.
\end{quote}

Now the talking is the room and the studying is the event. Neither version is more correct. They are about different things.
\end{grammarlens}

\subsection*{Your turn}
Say these aloud:

\begin{itemize}
\raggedright
  \item \textquotedblleft{}Estudiaba cuando habló Ana\textquotedblright{}
  \item the reversed version, and notice the story changed
\end{itemize}

\begin{tcolorbox}[breakable,colback=teal!4,colframe=teal!35!black,title={Before you move on}]
Which tense is the room? (The \textbf{imperfect}.) Which one walks into it? (The \textbf{preterite}.) Next: two verbs where the choice does not change the picture — it changes the \textbf{meaning}.
\end{tcolorbox}
